\section{El punto de partida de la negociación de 2027}

El salario mínimo busca proteger la remuneración del trabajo. Para que esa protección sea sostenible, su ajuste debe considerar tanto el costo de vida como la capacidad de producir de las empresas y las oportunidades de quienes todavía no acceden a un empleo formal. La inflación informa cuánto se ha encarecido el consumo. La productividad informa cuánto valor agregado real se produce por trabajador o por hora. Confundir ambas variables impide distinguir la compensación de precios de una mejora real del ingreso.

Este informe distingue los criterios para definir el nivel del mínimo y su actualización a partir de la revisión bibliográfica de Orozco y de casos documentados en América. Después examina la evolución del salario mínimo real y la productividad laboral desde 2022, con 2021 como base, y calcula escenarios para el cierre de 2026. El análisis colombiano utiliza los montos legales del salario básico, el IPC nacional y las series de productividad laboral del DANE. Las cifras anuales de productividad llegan a 2025 y las proyecciones de inflación se identifican por entidad y fecha.

La propuesta es que el incremento de 2027 compense la inflación de 2026 y no supere el 7\% bajo las proyecciones hoy disponibles. La razón es el punto de partida de la negociación. Entre 2021 y 2025, el mínimo ganó poder adquisitivo a una tasa muy superior al crecimiento de la productividad laboral. El aumento nominal de 23\% de 2026 añade una ganancia real proyectada de alrededor de 15\%. Repetir una mejora real importante exigiría justificar por qué la economía puede absorberla después de esos aumentos.

La discusión nacional ya contiene propuestas que superan la inflación esperada. El 27 de octubre, Bloomberg Línea reportó que Fenalco planteaba un límite de 9\% para 2027.\cite{debate} Con una inflación de referencia de 6,9\%, un ajuste nominal de 9\% excedería esa compensación en 2,1 puntos porcentuales. La diferencia equivale a una ganancia de 2,0\% al deflactar el ajuste con esa inflación. Por tanto, incluso una propuesta de un solo dígito puede incorporar una mejora real adicional.

El porcentaje de 7\% no constituye una regla legal. La fijación del mínimo colombiano considera varios criterios económicos y sociales, como recoge la motivación del Decreto 159 de 2026.\cite{decreto2026} La recomendación de este informe es una valoración económica para la concertación. Su fundamento combina la evolución acumulada del salario real, el crecimiento de la productividad y la incertidumbre sobre los efectos del aumento de este año.

\section{Qué debe determinar el salario mínimo}
\label{sec:criterios}

La revisión bibliográfica de Orozco organiza la discusión en tres decisiones. Primero, qué nivel del mínimo permite proteger los ingresos de los trabajadores. Segundo, cómo actualizarlo para compensar la inflación e incorporar mejoras reales. Tercero, qué condiciones de empleo, formalidad y capacidad empresarial deben considerarse al tomar esas decisiones.\cite{orozcoRevision} La distinción permite evaluar una propuesta de aumento sin suponer que el nivel inicial era adecuado ni que toda mejora real debe repetirse al año siguiente.

\subsection{Nivel, actualización y salvaguardas}

El \textbf{nivel de referencia} requiere examinar las necesidades de los trabajadores y sus familias y la posición del mínimo en la distribución salarial. Una canasta de consumo informa sobre los gastos que debe cubrir el hogar. La relación entre el mínimo y el salario mediano informa sobre la distancia frente al trabajador situado en el centro de la distribución. Son comparaciones diferentes. El presupuesto de una familia depende de su tamaño, de cuántos integrantes trabajan y de los precios locales. Por ello, un salario mínimo individual no puede compararse directamente con una línea de pobreza por persona sin precisar la composición del hogar. La OIT incorpora estas diferencias al analizar la suficiencia de los salarios.\cite{oitNecesidades}

La \textbf{actualización periódica} parte de un nivel ya fijado. Compensar la inflación evita que ese nivel pierda poder de compra por el aumento de los precios. Incorporar productividad permite que los trabajadores participen del mayor valor agregado que la economía produce por trabajador o por hora. Pero una regla de inflación más productividad no determina si el nivel inicial era adecuado. Si el mínimo era insuficiente, la regla no garantiza corregir esa insuficiencia. Si se había elevado mucho frente a la producción por trabajador, la regla tampoco evalúa cómo las empresas están absorbiendo ese cambio.

Las \textbf{salvaguardas} examinan la posibilidad de sostener el nivel y el ajuste propuestos. La proporción de trabajadores alcanzados por el mínimo, la contratación, las horas, la informalidad y el cumplimiento permiten observar dimensiones que no aparecen en la tasa de productividad nacional. La OIT advierte que la incidencia de un incremento depende también de cuántos trabajadores están afectados y de su efecto sobre la masa salarial.\cite{oitFactores} Estas variables deben interpretarse junto con la actividad económica, porque una caída de la contratación no identifica por sí sola una consecuencia del mínimo.

\subsection{Las experiencias de América responden a objetivos distintos}

El Cuadro \ref{tab:reglas} recoge seis mecanismos documentados en países de América. Los casos ilustran las modalidades discutidas por Orozco. No constituyen una clasificación exhaustiva del continente ni permiten establecer qué mecanismo funciona mejor por el número de países que lo adopta. Una regla publicada y el resultado de una negociación son objetos distintos de comparación.

\begin{table}[H]
\centering\small
\caption{Ejemplos de mecanismos de fijación y actualización del mínimo}
\label{tab:reglas}
\renewcommand{\arraystretch}{1.16}
\begin{tabularx}{\textwidth}{@{}>{\raggedright\arraybackslash}p{2.2cm}>{\raggedright\arraybackslash}X>{\raggedright\arraybackslash}X@{}}
\toprule
Caso & Mecanismo documentado & Distinción relevante\\
\midrule
Canadá, ámbito federal & Ajuste anual con el IPC del año calendario anterior. & Protege frente a precios. Su cobertura es federal y puede prevalecer un mínimo provincial superior.\cite{canada}\\[4pt]
Perú & Metodología de referencia que combina inflación esperada y productividad total de los factores. & La fórmula técnica debe distinguirse de la decisión y la oportunidad de los reajustes.\cite{peru}\\[4pt]
Costa Rica, sector privado & Metodología publicada con costo de vida, corrección por inflación y una porción del crecimiento real por habitante. & El PIB por habitante no equivale a productividad por trabajador. Incluye cláusulas de revisión.\cite{costarica}\\[4pt]
México, zona general & Monto Independiente de Recuperación más un factor de aumento. & Separa un componente de recuperación del ajuste porcentual. El esquema de 2026 difiere en la frontera norte.\cite{mexico}\\[4pt]
Colombia & Concertación con varios criterios económicos y sociales. & No existe una regla que imponga sumar mecánicamente inflación y productividad.\cite{decreto2026}\\[4pt]
Estados Unidos, ámbito federal & Modificación legislativa, sin indexación automática. & El mínimo federal es de USD 7,25 por hora desde 2009. Pueden aplicar pisos estatales superiores.\cite{eeuu}\\
\bottomrule
\end{tabularx}
\fuente{Elaboración del CJC a partir de la revisión de Orozco y las fuentes oficiales citadas. Son casos seleccionados. Para Perú se describe la metodología revisada en Moneda 199 (2024). Para Costa Rica se describe la metodología publicada por el Ministerio de Trabajo, no un cálculo del reajuste de 2027.}
\end{table}

Canadá muestra que la indexación con precios puede constituir una regla explícita de protección del salario. El mínimo federal se ajusta anualmente el primero de abril usando el IPC del año anterior. Si el piso provincial o territorial aplicable es superior, prevalece ese valor.\cite{canada} La experiencia permite identificar una modalidad de actualización, pero no demuestra que indexar solamente con inflación sea apropiado en cualquier país o para cualquier nivel inicial.

Perú y Costa Rica permiten distinguir las medidas que suelen agruparse bajo la palabra productividad. La metodología peruana revisada por el Banco Central de Reserva combina inflación esperada y variación de la productividad total de los factores.\cite{peru} La metodología costarricense publicada utiliza una parte del crecimiento del PIB real por habitante y contempla condiciones que llevan a revisar la aplicación automática de sus componentes.\cite{costarica} La PTF mide la parte del crecimiento que no se atribuye a los insumos productivos medidos. El PIB por habitante también depende de cuántas personas trabajan y de las horas que trabajan. Ninguno es intercambiable con el valor agregado por ocupado o por hora que utiliza este informe.

México incorpora una distinción adicional. La resolución para 2026 separa el Monto Independiente de Recuperación (MIR) del aumento porcentual ordinario. En la zona general se agregaron 17,01 pesos diarios al mínimo anterior y se aplicó un factor de 6,5\% sobre la suma. En la frontera norte el ajuste fue de 5\%, sin MIR.\cite{mexico} El caso muestra cómo una decisión de recuperación puede hacerse explícita y diferenciarse por ámbito. No prueba que Colombia enfrente el mismo rezago ni que deba aplicar el mismo porcentaje.

Estados Unidos ilustra el riesgo de depender de cambios legislativos sin una actualización automática. El mínimo federal sigue en USD 7,25 por hora desde julio de 2009. Con inflación positiva, mantener su valor nominal reduce el poder de compra. Sin embargo, varios estados tienen mínimos superiores, por lo que el piso federal no representa la remuneración mínima aplicable a todos los trabajadores estadounidenses.\cite{eeuu} La comparación relevante es entre mecanismos de actualización, sin equiparar ese congelamiento nominal prolongado con un ajuste colombiano que compense la inflación de un año.

\subsection{Qué implica la revisión para Colombia y para el ajuste de 2027}

La revisión comparada respalda una forma de ordenar la decisión, no un porcentaje único. La compensación de precios, la participación en mejoras productivas y la recuperación de un nivel considerado insuficiente persiguen objetivos diferentes. La concertación debería precisar qué parte del incremento propuesto corresponde a cada objetivo y con qué evidencia se justifica. El crecimiento del PIB, el PIB por habitante, la PTF y la productividad laboral no deben contarse como si fueran ganancias independientes que siempre pudieran sumarse.

Para 2027, este informe propone compensar la inflación de 2026 después de las ganancias reales acumuladas del mínimo. Es una recomendación para el próximo ajuste, no una regla permanente que excluya futuras mejoras reales. Tampoco certifica que el nivel actual cubra las necesidades de todos los hogares. Para resolver esa pregunta hacen falta indicadores de suficiencia y distribución salarial que conserven el tamaño del hogar, el número de perceptores de ingresos y la jornada laboral.

El fundamento del techo provisional de 7\% debe buscarse, por tanto, en el diagnóstico colombiano que se desarrolla a continuación. Los casos internacionales muestran que es posible separar los objetivos del ajuste y definir salvaguardas. No convierten el 7\% en un consenso internacional ni sustituyen el análisis de cuánto han crecido el salario real y la productividad en Colombia.


\section{El salario real y la productividad se separaron}

La Figura \ref{fig:historica} compara tres índices con base 2021 igual a 100. El primero mide el salario mínimo mensual dividido por el IPC promedio de cada año. Los otros dos encadenan las tasas de crecimiento de la productividad laboral que publica el DANE. La productividad por hora relaciona el valor agregado real con las horas trabajadas. La productividad por persona empleada lo relaciona con el número de trabajadores.

En 2025, el índice del salario mínimo real llegó a 113,5, el de productividad por hora a 101,2 y el de productividad por persona empleada a 95,4. El mínimo ganó 13,5\% de poder adquisitivo respecto de 2021. En el mismo periodo, la producción por hora aumentó 1,2\% y la producción por persona empleada cayó 4,6\%. La diferencia entre los crecimientos acumulados del salario real y la productividad por hora fue de 12,3 puntos porcentuales. Frente a la productividad por persona empleada fue de 18,2 puntos.\cite{ipc,productividad,salarios}

\begin{figure}[H]
\centering
\includegraphics[width=\textwidth]{salario_productividad_2021_2025.pdf}
\caption{Salario mínimo real y productividad laboral, 2021--2025}
\label{fig:historica}
\fuente{El salario se deflacta con IPC promedio anual. Productividad por hora y por persona empleada, enfoque de valor agregado. Índices construidos mediante encadenamiento de tasas anuales. 2024 provisional y 2025 preliminar para productividad. Las series no son medidas de productividad individual de los trabajadores que reciben el mínimo.}
\end{figure}

La separación ocurrió en varios años, pero no siguió una trayectoria uniforme. En 2022, el aumento nominal del mínimo no alcanzó a compensar la inflación del mismo año. El salario real promedio cayó 0,10\% y el salario real de diciembre cayó 2,70\%. En 2023, el salario real promedio aumentó 3,82\%, mientras la productividad por hora cayó 1,10\%. En 2024 y 2025, el salario real volvió a aumentar más que ambas medidas de productividad (Cuadro \ref{tab:anual}). La afirmación relevante es acumulada desde 2022, no una ganancia real del mismo tamaño en cada año.

\begin{table}[H]
\centering\small
\caption{Cambios anuales del salario mínimo y la productividad (\%)}
\label{tab:anual}
\begin{tabular}{lrrrrr}
\toprule
& Mínimo & \multicolumn{2}{c}{Mínimo real} & \multicolumn{2}{c}{Productividad}\\
\cmidrule(lr){3-4}\cmidrule(lr){5-6}
Año & nominal & Promedio & Diciembre & Por hora & Por persona\\
\midrule
2022 & 10,07 & -0,10 & -2,70 & -0,28 & -2,20 \\
2023 & 16,00 & 3,82 & 6,15 & -1,10 & -2,49 \\
2024 & 12,07 & 5,12 & 6,53 & 2,23 & 0,55 \\
2025 & 9,50 & 4,14 & 4,19 & 0,40 & -0,56 \\
\bottomrule
\end{tabular}
\fuente{Fuentes: DANE y montos legales. Cálculos del CJC. Cada variación compara el año con el anterior. Promedio y diciembre son dos deflactores distintos. Las diferencias mínimas con tasas oficiales de inflación pueden deberse al redondeo de los índices publicados.}
\end{table}

La diferencia entre las dos medidas de productividad también importa. Si el número de horas por persona cae, la producción por hora puede mejorar aunque la producción por trabajador se estanque. Por eso un aumento de la productividad por hora no debe trasladarse automáticamente al salario mensual sin examinar la jornada. La comparación del mínimo mensual con la productividad por persona aporta una referencia complementaria.

El diagnóstico tiene un alcance preciso. El mínimo es un piso legal y no la remuneración media. La productividad laboral describe el conjunto de la economía y no la productividad de quienes ganan el mínimo. La distancia entre los índices documenta que sus trayectorias se separaron. No mide directamente un sobrecosto de cada empresa, una pérdida de empleo ni la distancia frente a un salario óptimo.

\section{El aumento de 2026 cambia la base de la discusión}

El salario mínimo mensual de 2026 es \$1.750.905, frente a \$1.423.500 en 2025. El aumento del salario básico fue exactamente 23\%.\cite{decreto2026} Para calcular su ganancia real al final de 2026 hay que descontar la inflación del mismo periodo. Si los precios suben 6,9\%, el cálculo es
\[
\Delta SM^{real}_{2026}=\frac{1,23}{1,069}-1=15,1\%.
\]
Restar 6,9 a 23 produce una aproximación de 16,1 puntos porcentuales. La variación real exacta es menor porque se obtiene dividiendo los factores de crecimiento.

La inflación anual observada en septiembre de 2026 fue de 6,29\%, frente a 5,10\% al cierre de 2025. El dato de septiembre no es la inflación de todo 2026.\cite{ipc2026,ipc2025} Para evaluar el cierre deben utilizarse pronósticos identificables, sin confundirlos con resultados del DANE.

El informe de política monetaria de julio, publicado el 4 de agosto, proyectó una inflación de 6,9\% para diciembre de 2026. Las minutas de la reunión de septiembre reportaron expectativas de los analistas de 6,8\%. Bancolombia proyectó 6,7\% en su actualización del 23 de septiembre.\cite{ipm,minutas,bancolombia} La divulgación de la Encuesta de Opinión Financiera de octubre, publicada por Fedesarrollo el 21 de octubre, recoge una expectativa de 6,87\%.\cite{eof} Son pronósticos de fuentes y fechas diferentes, no los límites de un intervalo estadístico de confianza.

Con inflación de cierre entre 6,7\% y 6,9\%, la ganancia real del mínimo de 2026 estaría entre 15,1\% y 15,3\%. Incluso si la inflación llegara a 7\%, sería de 15,0\%. Esta ganancia proyectada supera ampliamente los crecimientos anuales recientes de la productividad laboral. No permite afirmar todavía cuánto la supera en el conjunto de 2026, porque el resultado anual de productividad no se ha publicado.

El escenario debe compararse con la historia usando el mismo deflactor. Con IPC de diciembre, el mínimo real aumentó 14,6\% entre 2021 y 2025, frente a 13,5\% con IPC promedio anual. Si la inflación de 2026 termina en 6,9\%, la ganancia acumulada entre diciembre de 2021 y diciembre de 2026 sería de aproximadamente 31,9\%. Esta última cifra combina un resultado histórico con un supuesto de cierre. No es un dato observado de 2026.

La trayectoria reciente de la economía también obliga a evitar conclusiones automáticas. El Banco de la República describe un mercado laboral con mayor ocupación y desempleo bajo durante 2026, al tiempo que identifica presiones de costos laborales y de demanda sobre los precios.\cite{ipm} Un buen resultado agregado de empleo no demuestra que el aumento del mínimo carezca de costos. Tampoco la mayor inflación demuestra que todo su incremento provenga del mínimo. Ambas interpretaciones requieren separar ese cambio de otros factores.

\section{Por qué ajustar con la inflación y limitar el aumento al 7\%}

Un ajuste cercano a la inflación de 2026 compensaría el encarecimiento de los bienes y servicios de este año y conservaría el punto de partida alcanzado por el salario. Después de la mejora real acumulada hasta 2025 y del salto de 2026, proponemos que la negociación de 2027 no añada otra mejora real importante. Con los pronósticos consultados, esa recomendación corresponde a un aumento cercano a 6,7\%--6,9\% y, en todo caso, no mayor a 7\%.

El Cuadro \ref{tab:escenarios} muestra las implicaciones en pesos. Si se adopta la referencia de 6,9\%, el salario básico subiría cerca de \$120.812 mensuales hasta \$1.871.717. Con el techo de 7\%, el incremento sería de \$122.563 hasta \$1.873.468. El auxilio de transporte requiere una decisión separada y no se incluye en estos cálculos.

\begin{table}[H]
\centering\small
\caption{Escenarios de inflación de 2026 y ajuste del salario básico de 2027}
\label{tab:escenarios}
\begin{tabular}{rrrr}
\toprule
Inflación 2026 & Ganancia real & Aumento mensual & Salario 2027\\
supuesto (\%) & del mínimo 2026 (\%) & si se indexa (pesos) & si se indexa (pesos)\\
\midrule
6,7 & 15,3 & 117.311 & 1.868.216 \\
6,8 & 15,2 & 119.062 & 1.869.967 \\
6,9 & 15,1 & 120.812 & 1.871.717 \\
7,0 & 15,0 & 122.563 & 1.873.468 \\
\bottomrule
\end{tabular}
\fuente{Cálculos del CJC. La segunda columna deflacta el aumento de 23\% de 2026. Las dos últimas aplican la inflación supuesta de 2026 al salario nominal de \$1.750.905 para ilustrar una indexación en 2027. Cifras redondeadas al peso. No incluyen auxilio de transporte.}
\end{table}

El techo de 7\% no se obtiene restando al salario las ganancias que superaron a la productividad. Ese procedimiento supondría que el nivel inicial era el adecuado, que la productividad agregada corresponde a los trabajadores afectados y que toda diferencia debe revertirse. El informe no demuestra esos supuestos. La recomendación es mantener las ganancias reales anteriores y moderar el siguiente ajuste, mientras se evalúa cómo las están absorbiendo las empresas y el mercado laboral.

El argumento a favor de una nueva mejora real es que algunos hogares siguen sin disponer de ingresos suficientes y que un salario más alto puede elevar el consumo de quienes lo reciben. Esa preocupación debe formar parte de la concertación. Pero la suficiencia del ingreso de los hogares no demuestra por sí sola que cada empleador pueda financiar el mismo aumento ni que el mínimo llegue a todos los hogares con menores ingresos. La OIT recomienda considerar conjuntamente las necesidades de los trabajadores, la productividad y el empleo.\cite{oit}

La literatura tampoco permite sostener que todo aumento del mínimo destruye empleo. Dube encuentra efectos generalmente pequeños sobre empleo y aumentos de remuneración en una revisión centrada en Estados Unidos, Reino Unido y otras economías desarrolladas.\cite{dube} Los estudios colombianos reunidos en ESPE 103 encuentran riesgos para la contratación formal y efectos sobre informalidad y precios, con diferencias según método y población.\cite{espe} Su lectura conjunta favorece evaluar el contexto colombiano y el tamaño del aumento, sin importar mecánicamente resultados de otro país ni tratar una simulación como un efecto observado.

La recomendación tiene una condición explícita. Si los pronósticos se mantienen por debajo de 7\%, compensar la inflación de 2026 y respetar ese techo son objetivos compatibles. Si la inflación observada o la mejor proyección disponible supera 7\%, habría que revisar el porcentaje. Mantener el límite en ese caso ya no sería una propuesta de compensación completa de precios. Asimismo, indexar con la inflación de 2026 no garantiza que el salario conserve su poder adquisitivo durante todo 2027, que dependerá de la inflación de ese nuevo año.

\section{La productividad y el acceso al empleo formal}

El estudio de Patiño, Chaparro y coautores sobre Antioquia permite distinguir dos costos del aumento del mínimo. Con registros de Comfama y Comfenalco, su modelo para unas 139.500 empresas privadas entre enero de 2022 y febrero de 2026 estima que un aumento real de 1\% del mínimo se asocia con aproximadamente 49.500 puestos formales menos durante un año. Cerca de 44.500 corresponden a empleos que dejan de crearse y 5.000 a destrucción de puestos. Alrededor del 90\% del efecto estimado ocurre, por tanto, por menor creación. El empleo puede crecer y, aun así, hacerlo menos que en ausencia del aumento. Son estimaciones del estudio para Antioquia, no un cálculo del impacto nacional ni del saldo atribuible al aumento de 23\% de 2026.\cite{chaparro}

El estudio también documenta compresión salarial. Entre quienes ganaban un mínimo en marzo de 2025, solo 4,7\% pasó al rango de uno a 1,5 mínimos en marzo de 2026, frente a 7,1\% en la transición anterior. Los salarios superiores crecieron menos que el piso, reduciendo la distancia entre remuneraciones. Pasar a un rango inferior medido en mínimos no implica necesariamente una caída del salario nominal. Esta evidencia descriptiva de movilidad salarial complementa la estimación sobre flujos de empleo.\cite{chaparro}

Para la concertación nacional, estos resultados muestran por qué observar únicamente despidos puede ocultar costos para quienes buscan trabajo. Después de las ganancias reales acumuladas, ajustar el mínimo con la inflación de 2026 evitaría añadir otra presión de origen regulatorio sobre las escalas salariales. Si la productividad aumenta mientras el mínimo real se mantiene, la distancia entre ambos puede reducirse. Así se daría margen para absorber el salto de 2026 y reducir el riesgo de nueva compresión. La indexación por sí sola no eleva la productividad ni restablece las diferencias salariales, que también dependen de la capacidad de las empresas para aumentar las remuneraciones superiores al mínimo.

El seguimiento de 2027 debería comparar contratación, permanencia, horas e ingresos en empresas con distinta proporción de trabajadores cerca del mínimo. También debería observar las entradas al empleo formal de jóvenes y personas con menor experiencia. Los resultados deben medirse frente a grupos y periodos comparables, porque el número de contratos por sí solo no distingue una respuesta al mínimo de una mejora general de la actividad económica.

El borrador de Orozco sobre salario mínimo y empleo juvenil plantea estudiar mecanismos para facilitar el ingreso de los jóvenes a la formalidad.\cite{orozcoBoletin} Un mínimo diferenciado por edad constituye una reforma distinta del porcentaje general para 2027. Requiere comparar sus efectos con subsidios a la contratación y formación, además de examinar sustitución de trabajadores y rotación. La mayor concentración de jóvenes en salarios bajos no prueba por sí sola que reducirles el piso salarial mejore su trayectoria laboral.

\section{Conclusión}

El mínimo ganó 13,5\% de poder adquisitivo promedio entre 2021 y 2025, mientras la productividad por hora aumentó 1,2\% y la productividad por persona empleada cayó 4,6\%. El aumento nominal de 23\% en 2026 añade una ganancia real proyectada cercana a 15\%. La evidencia disponible sustenta moderar el siguiente ajuste, aunque no permite calcular un porcentaje óptimo ni una brecha definitiva de productividad para 2026.

La revisión bibliográfica permite separar la suficiencia del nivel, la compensación de precios y las condiciones para una mejora real. Los mecanismos internacionales utilizan combinaciones distintas de esos criterios y no establecen un porcentaje universal. Con base en el diagnóstico colombiano, proponemos que el mínimo de 2027 aumente cerca de la inflación de 2026 y no supere el 7\% con las proyecciones actuales. La referencia del equipo técnico del Banco de la República es 6,9\%. La cifra deberá actualizarse antes de la decisión. La propuesta conserva las ganancias reales anteriores sin establecer un congelamiento real permanente.
\clearpage
